\chapter{相对论量子力学与原子精细结构}
\label{ch:relativity}

\section{精细结构与Dirac理论}
原子两能级之间的相干保持得越久，相应跃迁的谱线就可以越窄。若两条谱线的间隔超过线宽，仪器分辨率也足够，就能分别测出两个共振频率。
改变外磁场，观察哪些谱线在零场仍然分开。要解释这些频率间隔，就要找出哪些相互作用使\chapref{ch:hydrogen}中原本简并的能级发生了不同的能移。改变光的偏振和原子的初态，有助于辨认每条线对应的跃迁。

谱学中 $n$ 标记主能级，$\ell$ 标记轨道角动量，$j$ 标记轨道角动量与电子自旋合成的总角动量。
例如氢的 $2\orbitalP_{1/2}$ 指 $n = 2,\ell = 1,j = 1/2$。若同时写出电子组态和光谱项，同一能级记作
$2\orbitalP\,{}^2\termP_{1/2}$，其中左上角的 $2$ 是自旋多重度。
相对论修正能有多大？在\chapref{ch:hydrogen}的点核模型中再取核质量无穷大的极限，
令 $m$ 为电子质量，$\alpha = e^2/(4\pi\varepsilon_0\hbar c)$。在低主量子数的氢样原子中，
特征动量 $p \sim \hbar/a_Z = mcZ\alpha$。展开自由粒子的相对论能量，得到
\begin{equation}
    \sqrt{m^2c^4 + p^2c^2}
    = mc^2 + \frac{p^2}{2m} - \frac{p^4}{8m^3c^2} + \cdots .
\end{equation}
非相对论动能为 $mc^2(Z\alpha)^2$ 量级，第一项相对论修正为 $mc^2(Z\alpha)^4$ 量级。
对氢取 $Z = 1$，修正与非相对论动能之比为 $\alpha^2$ 量级，其中 $\alpha^2 \simeq 5 \times 10^{-5}$。
因此，能分辨主能级间隔的仪器未必能分辨精细结构。理论要达到相应精度，还需保留与相对论动能修正同阶的自旋和静电势效应。

例如，非相对论 Coulomb 模型预言 $2\orbitalP_{1/2}$ 和 $2\orbitalP_{3/2}$ 能量相同。Dirac 理论把电子自旋纳入相对论动力学，由此得到这两条能级间的精细结构(fine structure)分裂。
但点核 Dirac 模型仍预言 $2\orbitalS_{1/2}$ 与 $2\orbitalP_{1/2}$ 简并。Lamb--Retherford 实验测出了这两态的能量差\cite{lamb}；解释这一 Lamb 位移还需要辐射场的量子理论。

自由相对论能量满足$E^2 = c^2p^2 + m^2c^4$。
能量为 $E$ 的定态带有时间因子 $\exp(-\i Et/\hbar)$。对这个因子求两次时间导数就会带出 $-E^2/\hbar^2$，按能量平方关系写出的波动方程含时间二阶导数。
Dirac选择让态仍按时间的一阶方程演化，同时让空间方程也只含一阶导数。
在位置表象里一阶空间导数对应一个动量因子，这就提示尝试在动量中线性的自伴Hamilton算符
\begin{equation}
    \op{H}_{\diracLabel} = c\opbm{\alpha} \cdot \opbm{P} + \op{\beta} mc^2.
\end{equation}

取 $\op{\alpha}_x,\op{\alpha}_y,\op{\alpha}_z,\op{\beta}$ 为只作用在内部空间上的算符，因而与空间动量 $\op{P}_i$ 对易。
粗体 $\opbm{\alpha}$ 的三个分量是 $\op{\alpha}_i$，标量 $\alpha$ 仍表示精细结构常数。
将Hamilton算符平方，
\[
    \op{H}_{\diracLabel}^2 =
    \frac{c^2}{2}\sum_{ij}\{\op{\alpha}_i,\op{\alpha}_j\}\op{P}_i\op{P}_j
    + mc^3\sum_i\{\op{\alpha}_i,\op{\beta}\}\op{P}_i
    + \op{\beta}^2m^2c^4.
\]
要求此式对任意动量都等于$(c^2\opbm{P}^2 + m^2c^4)\one$，
就分别固定二次项、一次项和常数项的系数。

\begin{definition}{Hamilton形式的Dirac矩阵}{qm-dirac-matrices}
    三维空间的 Dirac Hamilton 量中，内部算符 $\op{\alpha}_x,\op{\alpha}_y,\op{\alpha}_z,\op{\beta}$ 自伴且无量纲。它们在内部正交归一基中的表示称为 Dirac 矩阵。这些算符满足
    \begin{subequations}
        \begin{align}
            \{\op{\alpha}_i,\op{\alpha}_j\} &= 2\delta_{ij}\one,\\
            \{\op{\alpha}_i,\op{\beta}\} &= 0,\\
            \op{\beta}^2 &= \one.
        \end{align}
    \end{subequations}
\end{definition}
为什么自旋的二分量空间还不够？
若内部空间是二维的，三个相互反对易且平方为恒等的 $\op{\alpha}_i$ 可取为三个 Pauli 算符（至多相差基底和轴向的选择）。
再写$\op{\beta} = b_0\one + \bm b \cdot \opbm{\sigma}$，
要求$\{\op{\beta},\op{\sigma}_i\} = 0$给出
$2b_0\op{\sigma}_i + 2b_i\one = 0$。
对三个$i$同时成立只可能有$b_0 = b_x = b_y = b_z = 0$，
这与$\op{\beta}^2 = \one$冲突。
因此在三维空间需要更大的内部表示。三分量也无解：$\op{\beta}$ 的本征值只能是 $+1$ 或 $-1$，
而可逆的 $\op{\alpha}_x$ 把 $\op{\beta}$ 的正本征子空间映到负子空间，反之亦然，两者必须等维。

在三维空间，内部空间的最小复维数为四。选取 $\op{\beta}$ 对角的 Dirac 基，每一块内依次取自旋基 $\ket{+z},\ket{-z}$。
以 $\sigma_i$ 表示 $\op{\sigma}_i$ 在这组自旋基中的矩阵，$I_2$ 表示二阶单位矩阵，可取
\begin{equation}
    \op{\alpha}_i \longleftrightarrow \begin{pmatrix}0&\sigma_i\\\sigma_i&0\end{pmatrix},\qquad
    \op{\beta} \longleftrightarrow \begin{pmatrix}I_2&0\\0&-I_2\end{pmatrix}.
\end{equation}
四分量态 $\ket{\Psi} \in \HH_{\spatial} \otimes \CC^4$ 遵循 Dirac 方程。对动量大小为 $p$ 的自由粒子，$\op H_{\diracLabel}$ 的平方给出两支能量的大小，
\begin{subequations}
    \begin{align}
        \i\hbar\,\dv{\ket{\Psi}}{t} &= \op H_{\diracLabel}\ket{\Psi},\\
        E_\pm(p) &= \pm\sqrt{m^2c^4 + c^2p^2}.
    \end{align}
\end{subequations}
在静止系中 $\op{H}_{\diracLabel} = \op{\beta} mc^2$，上面两个分量属于 $+mc^2$，下面两个属于 $-mc^2$，
每一支都容纳两种自旋投影。一旦有动量，$\opbm{\alpha} \cdot \opbm{P}$ 的非对角块就把上下分量联系起来。
负能分支标志着相对论单粒子描述的边界；涉及粒子与反粒子成对产生的过程需要量子场论。

\section{Pauli近似}
Dirac 方程用四个分量描述自由电子，弱磁场中的自旋进动则可以用二分量态描述。这两种描述如何衔接？对同一个低速电子，四分量方程应能给出两种自旋投影以及磁场引起的能量分裂。
为此先对能量接近 $+mc^2$ 的一支求低速极限。

经典外场 $\bm A(\bm x,t)$、$\phi(\bm x,t)$ 和 $\bm B(\bm x,t)$ 都是给定的数值函数。按\chapref{ch:em}的最小耦合，用它们在位置读数 $\bm x$ 上的值构造乘法算符
\[
    \opbm{A} = \bm A(\opbm{X},t),\qquad
    \op{\phi} = \phi(\opbm{X},t),\qquad
    \opbm{B} = \bm B(\opbm{X},t).
\]
例如对位置 $\bm x$ 和内部基标号 $s$ 取分量，有 $\mel{\bm x,s}{\op\phi}{\Psi(t)} = \phi(\bm x,t)\braket{\bm x,s}{\Psi(t)}$。
用机械动量 $\opbm{\pi} = \opbm{P} - q\opbm{A}$ 写出外场中的 Dirac 算符，
\begin{equation}
    \op{H}_{\diracLabel} = c\opbm{\alpha} \cdot \opbm{\pi} + \op{\beta} mc^2 + q\op{\phi}.
\end{equation}

提出静能相位，在上述 Dirac 基中用两个二分量右矢表示四分量态，
\begin{equation}
    \ket{\Psi(t)}
    \longleftrightarrow \e^{-\i mc^2t/\hbar}
    \begin{pmatrix}\ket{\varphi(t)}\\\ket{\chi(t)}\end{pmatrix}.
\end{equation}
$\ket{\varphi},\ket{\chi}$各自都属于
$\HH_{\spatial} \otimes \CC^2$；
它们是同一个态的两个块，不分别要求单位归一。
正能低速状态中，下块$\ket{\chi}$相对上块小一个$v/c$量级。

去掉静能相位并将态分成两个二分量块$\ket{\varphi},\ket{\chi}$，得
\begin{subequations}
    \begin{alignat}{1}
        \i\hbar\,\dv{\ket{\varphi}}{t} &= q\op{\phi}\ket{\varphi} + c\opbm{\sigma} \cdot \opbm{\pi}\ket{\chi},\\
        \i\hbar\,\dv{\ket{\chi}}{t} &= c\opbm{\sigma} \cdot \opbm{\pi}\ket{\varphi} + (q\op{\phi} - 2mc^2)\ket{\chi}.
    \end{alignat}
\end{subequations}
要消去下块，先在第二式中把含 $\ket\chi$ 的项移到同一侧，
\[
    \biggp{2mc^2 + \i\hbar\,\dd t - q\op{\phi}}\ket{\chi}
        = c\,\opbm{\sigma} \cdot \opbm{\pi}\ket{\varphi}.
\]
提出静能相位后，若剩余能量、场致能量和 $\hbar$ 乘以相关变化频率都远小于 $2mc^2$，
且负能支混合很弱，左边的主项就是 $2mc^2\ket{\chi}$。
于是小分量可由大分量近似算出，
\begin{equation}
    \ket{\chi} \simeq \frac{\opbm{\sigma} \cdot \opbm{\pi}}{2mc}\ket{\varphi}.
    \label{eq:dirac-small-component}
\end{equation}
在机械动量集中于 $\bm\pi_0$ 的波包中，小分量与大分量的范数之比约为 $\abs{\bm\pi_0}/(2mc)$。
由Pauli代数与式\eqref{eq:kincomm}，
\begin{equation}
    (\opbm{\sigma} \cdot \opbm{\pi})^2 = \opbm{\pi}^2 - q\hbar\opbm{\sigma} \cdot \opbm{B}.
\end{equation}

逐分量计算为
\[
    \begin{alignedat}{1}
        (\opbm{\sigma} \cdot \opbm{\pi})^2 &= \sum_{ij}(\delta_{ij}\one + \i\epsilon_{ijk}\op{\sigma}_k)\op{\pi}_i\op{\pi}_j\\
        &= \opbm{\pi}^2 + \frac{\i}{2}\epsilon_{ijk}\op{\sigma}_k[\op{\pi}_i,\op{\pi}_j]\\
        &= \opbm{\pi}^2 - \frac{q\hbar}{2}
            \epsilon_{ijk}\epsilon_{ijl}\op{\sigma}_k\op{B}_l
            = \opbm{\pi}^2 - q\hbar\opbm{\sigma} \cdot \opbm{B}.
    \end{alignedat}
\]
第二行使用 $\epsilon_{ijk}$ 的反对称性。最后一行对重复的 $i,j$ 求和，用到
\[
    \sum_{i,j}\epsilon_{ijk}\epsilon_{ijl} = 2\delta_{kl}.
\]

将式\eqref{eq:dirac-small-component}代回上块方程，再用上述恒等式，便得到正能低速态的 Pauli 算符，
\begin{equation}
    \op{H}_{\mathrm{P}} = \frac{\opbm{\pi}^2}{2m} + q\op{\phi} - \frac{q\hbar}{2m}\opbm{\sigma} \cdot \opbm{B}.
\end{equation}
将最后一项写成磁矩的势能 $-\opbm\mu \cdot \opbm B$。代入 $\opbm S = \hbar\opbm\sigma/2$ 可知 $\opbm\mu = (q/m)\opbm S$。按 $\opbm\mu = g_{\spin}[q/(2m)]\opbm S$ 定义自旋 $g$ 因子，便得到 $g_{\spin} = 2$。
辐射修正使电子的 $g_{\spin}$ 略微偏离这个 Dirac 理论值，多出的磁矩称为反常磁矩(anomalous magnetic moment)。

\sectionstar{Foldy--Wouthuysen变换}
\label{sec:fw}
上节保留的最低阶已经给出电子的自旋磁矩。现在撤去外磁场，只让电子处在原子核的静电势中。
要看清零外场的精细结构，还须保留相对静能为 $(v/c)^4$ 量级的能量修正。用幺正变换逐阶减小上下两块之间的耦合，
可以把这一阶对正能态有贡献的项留在同一个二分量 Hamilton 算符里。

\begin{definition}{相对于 \texorpdfstring{$\op{\beta}$}{β} 的奇、偶算符}{qm-dirac-blocks}
    与 $\op{\beta}$ 对易的算符称为偶算符(even operator)，它将 $\op{\beta}$ 的每个本征子空间映回自身。
    与 $\op{\beta}$ 反对易的算符称为奇算符(odd operator)，它将一个本征子空间映入另一个。
    在 $\op{\beta}$ 对角的表示中，前者只有对角块，后者只有非对角块。
\end{definition}
记静电势能算符为 $\op{V} = V(\opbm{X})$，则 $\op{H}_{\diracLabel} = \op{\beta} mc^2 + \op{O} + \op{V}$，其中 $\op{O} = c\opbm{\alpha} \cdot \opbm{P}$。
势能 $\op{V}$ 满足 $[\op{\beta},\op{V}] = 0$，动量项 $\op{O}$ 则因 Dirac 矩阵的反对易关系满足 $\{\op{\beta},\op{O}\} = 0$。

可先把这个变换看作两能级对角化的推广。对只沿 $z$ 运动的自由粒子，取 $\op{\sigma}_z$ 的 $+1$ 本征子空间并固定动量 $p_z$ 后，
上下两块缩成一个 $2 \times 2$ 矩阵，对角元为 $\pm mc^2$，非对角元为 $cp_z$。
当 $\abs{p_z} \ll mc$，两支的间隔 $2mc^2$ 远大于耦合，小角度的基底旋转就能消去最低阶耦合。
同样可以选择一个生成元，使它与静能项的对易子抵消 $\op{O}$。

选择无量纲反 Hermite 生成元
$\op{S} = \op{\beta}\op{O}/(2mc^2)$。在动量表象中，每个固定动量 $\bm p$ 对应一个反 Hermite 矩阵 $S(\bm p)$。对它取矩阵指数，在 $S(\bm p)$ 的各本征空间乘以相应的指数，就得到幺正矩阵。以这些矩阵逐点作用于动量振幅，定义全空间上的幺正变换 $\op U_{\FW} = \e^{\op S}$。
下式按嵌套对易子作低速展开，计算只保留所需阶次。涉及无界算符的对易子先在有关乘积均有定义的测试波包上计算。
\begin{equation}
    \op{U}_{\FW}\op{H}_{\diracLabel}\op{U}_{\FW}\adjoint
    = \op{H}_{\diracLabel} + [\op{S},\op{H}_{\diracLabel}]
    + \frac12[\op{S},[\op{S},\op{H}_{\diracLabel}]] + \cdots.
\end{equation}
由于$[\op{S},\op{\beta} mc^2] = -\op{O}$，
最低阶奇项被消去。计算所需的后续对易子为
\begin{subequations}
    \begin{alignat}{1}
        [\op{S},\op{O}] &= \frac{\op{\beta}\op{O}^2}{mc^2},\\
        [\op{S},[\op{S},\op{O}]] &= -\frac{\op{O}^3}{m^2c^4},\\
        [\op{S},[\op{S},[\op{S},\op{O}]]] &= -\frac{\op{\beta}\op{O}^4}{m^3c^6},\\
        [\op{S},[\op{S},\op{V}]] &= -\frac{[\op{O},[\op{O},\op{V}]]}{4m^2c^4}.
    \end{alignat}
\end{subequations}
例如第一式中$\op{O}\op{\beta} = -\op{\beta}\op{O}$，
所以$[\op{\beta}\op{O},\op{O}] = 2\op{\beta}\op{O}^2$。
用 $[\op{S},\op{\beta}mc^2] = -\op{O}$ 合并来自动量项和静能项的贡献，二阶偶项为
\begin{equation}
    [\op{S},\op{O}] + \frac12[\op{S},[\op{S},\op{\beta}mc^2]]
        = \frac12[\op{S},\op{O}] = \frac{\op{\beta}\op{O}^2}{2mc^2}.
\end{equation}
同理，四阶偶项来自 $\op O$ 的三重对易子和 $\op{\beta}mc^2$ 的四重对易子，
\begin{equation}
    \frac1{3!}[\op{S},[\op{S},[\op{S},\op{O}]]]
    + \frac1{4!}[\op{S},[\op{S},[\op{S},
        [\op{S},\op{\beta} mc^2]]]]
    = -\frac{\op{\beta}\op{O}^4}{8m^3c^6}.
\end{equation}
因为$1/3! - 1/4! = 1/8$，而所列三重对易子本身带负号。
势能的二重对易子直接乘展开系数$1/2$，
给出$-[\op{O},[\op{O},\op{V}]]/(8m^2c^4)$。

在原子低速态中，$\op{O}$ 的典型矩阵元为 $mc^2(v/c)$ 量级，$\op{V}$ 的典型矩阵元为 $mc^2(v/c)^2$ 量级。按 $v/c$ 计阶，
剩余奇项从三阶开始。用同样方法再变换一次消去它，
到四阶的偶部分保持为\cite{littlejohnfw}
\begin{equation}
    \op{H}_{\FW}
    = \op{\beta}\Bigp{mc^2 + \frac{\op{O}^2}{2mc^2}
        - \frac{\op{O}^4}{8m^3c^6}}
    + \op{V} - \frac{[\op{O},[\op{O},\op{V}]]}{8m^2c^4}.
    \label{eq:fw-even}
\end{equation}
对于静电势，
$[\op{O},\op{V}] = -\i\hbar c\opbm{\alpha} \cdot (\grad V)(\opbm{X})$。
再用$\op{\alpha}_i\op{\alpha}_j = \delta_{ij}\one + \i\epsilon_{ijk}\op{\Sigma}_k$，
其中 $\op{\Sigma}_k$ 在 Dirac 基中的矩阵为 $\diag(\sigma_k,\sigma_k)$，
得到
\begin{equation}
    [\op{O},[\op{O},\op{V}]]
    = -\hbar^2c^2(\laplacian V)(\opbm{X})
    - 2\hbar c^2\opbm{\Sigma} \cdot \bigb{(\grad V)(\opbm{X}) \times \opbm{P}}.
\end{equation}
这里$\curl\grad V = 0$使反对称二阶导数项消失。
取正能上块，$\op{\beta}$ 限制为 $\one_2$，$\opbm{\Sigma}$ 限制为 $\opbm{\sigma}$。
沿用\chapref{ch:hydrogen}的距离算符 $\op R = \abs{\opbm X}$ 和方向算符 $\opbm n$。中心势的梯度为 $\grad V = V'(r)\bm n$，对应的算符为 $(\grad V)(\opbm X) = V'(\op R)\opbm n$。因 $\opbm L = \opbm X \times \opbm P$，有
\[
    (\grad V)(\opbm{X}) \times \opbm{P}
        = \frac{V'(\op{R})}{\op{R}}\opbm{L}
\]
其中 $V'(\op R)/\op R$ 按距离算符的谱定义为 $g(\op R)$，径向因子为 $g(r) = V'(r)/r$。再代入 $\opbm S = \hbar\opbm\sigma/2$，便得到下一节的有效 Hamilton 量。

\section{氢原子的精细结构}
氢的 $2\orbitalP_{1/2}$ 与 $2\orbitalP_{3/2}$ 具有相同的主量子数和轨道角动量，
在非相对论的 Coulomb 势中能量相同。加入相对论修正后，这两条能级还会重合吗？按\chapref{ch:approximations}的简并微扰理论，要在固定主量子数 $n$ 的简并子空间中求修正矩阵的本征值。

对固定点核和静电中心势 $V(r)$，保留到 $1/c^2$ 阶，正能支的有效 Hamilton 算符为
\begin{equation}
    \op{H}_\eff = mc^2\one + \frac{\opbm{P}^2}{2m} + V(\op{R})
            - \frac{\opbm{P}^4}{8m^3c^2}
            + \frac{1}{2m^2c^2}\frac{V'(\op{R})}{\op{R}}\opbm{L} \cdot \opbm{S}
            + \frac{\hbar^2}{8m^2c^2}(\laplacian V)(\opbm{X}).
    \label{eq:paulicorrections}
\end{equation}
后三项依次为相对论动能、自旋轨道修正(spin--orbit correction)和 Darwin 修正。
自旋轨道项来自式\eqref{eq:fw-even}的双对易子，Darwin项来自其中的$\laplacian V$。
二者与相对论动能项同阶，因而需要一起计算\cite{mitfine}。

动能修正和 Darwin 项在空间旋转下不变，对自旋都按恒等算符作用。自旋轨道项的径向因子也与角动量对易，
而 $\opbm L \cdot \opbm S = (\opbm J^2 - \opbm L^2 - \opbm S^2)/2$。所以三个修正都保持 $\ell,j,m_j$。
在固定 $n$、电子自旋 $s = 1/2$ 的子空间里，每组允许的 $\ell,j,m_j$ 只对应一个耦合基矢 $\ket{\psi_0} = \ket{n,\ell,s;j,m_j}$。
修正矩阵在这组基中已经对角，所求一级能移就是三个修正项在同一基矢上的期望值之和。若分别用轨道和自旋的投影 $m_\ell,m_s$ 作标签，自旋轨道项一般会混合总投影相同的乘积态。

在点核近似下，为什么电子的 Darwin 能移只取决于它在核处的概率密度？取原子核为原点，用 $\bm r$ 表示从核指向电子的位置矢量，$r = \abs{\bm r}$ 为电子到核的距离。电子的 Coulomb 势能为 $V(r) = -\kappa/r$，在 $r > 0$ 有 $\laplacian V = 0$，
而 $\grad V = (\kappa/r^2)\bm n$ 穿过任意以原点为心的球面的通量都是 $4\pi\kappa$。
由散度定理，$\laplacian V$ 在含原点的小球内的体积分也总是 $4\pi\kappa$。
因此 $\laplacian V$ 在分布意义下集中于原点。对本节归一化的 Coulomb 束缚本征态，选定自旋量子化轴，取连续的位置波函数 $\psi_{m_s}(\bm r) = \braket{\bm r,m_s}{\psi_0}$，其中 $m_s = \pm 1/2$。
连续性使原点值由邻近的振幅唯一确定。一般平方可积波函数可以只改动原点的值而仍表示同一个态。
下面的 $\ave{(\laplacian V)(\opbm X)}$ 表示分布作用于束缚态的位置密度所给的二次型。Darwin 项对自旋按恒等算符作用，因而两个自旋分量在原点的概率密度相加，
\begin{subequations}
    \begin{align}
        \laplacian V(\bm r) &= 4\pi\kappa\delta^{(3)}(\bm r),\\
        \ave{(\laplacian V)(\opbm X)} &= 4\pi\kappa\sum_{m_s = \pm 1/2}\abs{\psi_{m_s}(\bm 0)}^2 .
    \end{align}
\end{subequations}
本节采用固定点核，令$m = m_\electron$、$\kappa = Z\alpha\hbar c$；有限核质量的反冲修正另行加入。
计算所需的Coulomb矩阵元为
\begin{subequations}
    \begin{alignat}{1}
        \ave{\op{R}^{-2}} &= \frac1{a_Z^2n^3(\ell + 1/2)},\\
        \ave{\op{R}^{-3}} &= \frac1{a_Z^3n^3\ell(\ell + 1/2)(\ell + 1)}\quad(\ell > 0),\\
        \ave{\delta^{(3)}(\opbm{X})} &= \frac{\delta_{\ell0}}{\pi a_Z^3n^3},\qquad
            a_Z = \frac{\hbar^2}{m\kappa}.
    \end{alignat}
\end{subequations}
前两式是用标量权重 $r^{-2}$、$r^{-3}$ 对径向概率求均值，第三式取原点密度。它们是标准 Coulomb 本征态的结果\cite{mitfine}，
它们与\chapref{ch:hydrogen}的能谱一起构成精细结构计算的零级数据。
对于$\ell = 0$，轨道角动量算符直接给出零，自旋轨道项随之消失，而Darwin项有有限贡献。此时 $\ave{\op{R}^{-3}}$ 发散。

为核对合并步骤，令$C = mc^2(Z\alpha)^4/(2n^3)$。
用 $\opbm{P}^2\ket{\psi_0} = 2m(E_n + \kappa/\op{R})\ket{\psi_0}$ 计算 $\ave{\opbm{P}^4} = \|\opbm{P}^2\ket{\psi_0}\|^2$，得

\begin{equation}
    \begin{alignedat}{1}
        \ave{\opbm{P}^4} &= 4m^2\bigl\|(E_n + \kappa/\op{R})\ket{\psi_0}\bigr\|^2\\
        &= 4m^2\Bigb{
            E_n^2 + 2E_n\kappa\ave{\op{R}^{-1}}
            + \kappa^2\ave{\op{R}^{-2}}}.
    \end{alignedat}
\end{equation}
由$\kappa\ave{\op{R}^{-1}} = -2E_n$，
方括号变成$-3E_n^2 + \kappa^2\ave{\op{R}^{-2}}$。
自旋轨道项可从合成角动量 $\opbm J = \opbm L + \opbm S$ 求出。对电子取 $s = 1/2$，有
\begin{equation}\label{eq:qm-spin-orbit-expectation}
    \ave{\opbm L \cdot \opbm S}
    = \frac12\ave{\opbm J^2 - \opbm L^2 - \opbm S^2}
    = \frac{\hbar^2}{2}\bigb{j(j + 1) - \ell(\ell + 1) - 3/4}.
\end{equation}
将各个期望值代入有效 Hamilton 算符，得到三个能量修正，

\begin{subequations}
    \begin{alignat}{1}
        \Delta E_{\kinetic} &= C\biggp{\frac3{4n} - \frac1{\ell + 1/2}},\\
        \Delta E_{\spinorbit} &= C\frac{j(j + 1) - \ell(\ell + 1) - 3/4}
            {2\ell(\ell + 1/2)(\ell + 1)}\quad(\ell > 0),\\
        \Delta E_{\mathrm{D}} &= C\delta_{\ell0}.
    \end{alignat}
\end{subequations}
对$j = \ell \pm 1/2$分别化简，或对$\ell = 0,j = 1/2$加Darwin项，统一得到
\begin{equation}
    \Delta E_{nj} = \frac{mc^2(Z\alpha)^4}{2n^3}
    \biggp{\frac3{4n} - \frac1{j + 1/2}}.
    \label{eq:fine}
\end{equation}

对 $\ell > 0$，自旋轨道项的分子 $j(j + 1) - \ell(\ell + 1) - 3/4$ 在两支上分别为 $\ell$ 和 $-(\ell + 1)$。因此括号内与 $\ell$ 有关的两项满足
\begin{alignat}{2}
    j &= \ell + 1/2,&\qquad
    -\frac1{\ell + 1/2} + \frac1{2(\ell + 1/2)(\ell + 1)}
        &= -\frac1{\ell + 1} = -\frac1{j + 1/2},\\
    j &= \ell - 1/2,&
    -\frac1{\ell + 1/2} - \frac1{2\ell(\ell + 1/2)}
        &= -\frac1{\ell} = -\frac1{j + 1/2}.
\end{alignat}
$\ell = 0$时动能项中的$-2$加上Darwin项的$+1$，
也给出$-1/(j + 1/2) = -1$。

修正只依赖 $n,j$，因此两个 $2\orbitalP$ 态按 $j$ 分裂，而 $2\orbitalS_{1/2}$ 与 $2\orbitalP_{1/2}$ 在此阶仍简并（图\ref{fig:hydrogen-fine-lamb}）。

\begin{figure}[htbp]
    \centering\tikzfigure{hydrogen-fine-lamb}
    \caption{零外场氢原子 $n = 2$ 能级，忽略超精细结构。中列按式\eqref{eq:fine}画出相对非相对论能级的下移；虚线连接各模型中相互对应的能级。右列以 $2\orbitalP_{1/2}$ 为参照，放大两态间的 Lamb 间隔 $h\nu_{\mathrm{L}} = E(2\orbitalS_{1/2}) - E(2\orbitalP_{1/2})$。}
    \label{fig:hydrogen-fine-lamb}
\end{figure}

上述低速展开要求 $Z\alpha \ll 1$。当 $Z\alpha$ 不再很小时，需要直接求解 Dirac 方程。要与高 $Z$ 离子的精密光谱比较，还须计入有限核大小及更高阶修正。

\begin{example}{氢原子的精细结构能移}{qm-hydrogen-fine-terms}
    哪些同阶修正共同决定氢原子的精细结构？在无限重的点核近似下，合并后的能量依赖哪些量子数？

    \par\smallskip\textbf{解}\quad 相对论动能、自旋轨道和 Darwin 三项合并后，$\ell$ 的单独依赖消去，留下式\eqref{eq:fine}中的 $n,j$ 依赖。
    有限核质量还引入反冲修正，应与约化质量效应共同计算。
\end{example}

\section{超精细结构与Lamb位移}
更高分辨率的氢谱显示，核磁矩的耦合使同一电子能级分裂。扣除这层分裂后，电子能级之间仍有 Dirac 理论未能解释的频移。

核自旋算符 $\opbm I$ 和电子角动量算符 $\opbm J$ 耦合成总角动量 $\opbm F = \opbm I + \opbm J$。沿用上节的 $j$ 表示电子角动量量子数，核自旋量子数记为 $I$，故 $F = \abs{I - j},\ldots,I + j$。在各向同性的超精细相互作用(hyperfine interaction)模型中，实常量 $A$ 具有能量量纲。将 $\opbm F^2$ 展开，可将相互作用写成
\begin{equation}
    \op H_{\hyperfine} = \frac{A}{\hbar^2}\opbm I \cdot \opbm J
    = \frac{A}{2\hbar^2}\bigp{\opbm F^2 - \opbm I^2 - \opbm J^2}.
\end{equation}
在共同本征态 $\ket{I,j;F,m_F}$ 中，三个平方算符分别取本征值 $\hbar^2 F(F + 1)$、$\hbar^2 I(I + 1)$ 和 $\hbar^2 j(j + 1)$，因此能移为
\begin{equation}
    E_F = \frac{A}{2}\bigb{F(F + 1) - I(I + 1) - j(j + 1)}.
\end{equation}
上述模型描述磁偶极超精细耦合。电四极相互作用要在给定的 $j$ 能级内产生一阶能移，还需要核自旋 $I \geq 1$ 且 $j \geq 1$。这是秩为 $2$ 的球张量在核与电子的角动量空间中分别满足三角条件的结果。
核电荷的有限分布改变静电势，多电子关联来自电子之间的相互作用；这两种能移应与核磁矩耦合引起的劈裂分别计算。

\subsection{Lamb--Retherford实验}
\label{ss:lamb-experiment}

1947 年，Lamb 与 Retherford 制备了一束处于亚稳态 $2\orbitalS_{1/2}$ 的氢原子。$2\orbitalS \to 1\orbitalS$ 的单光子电偶极跃迁受宇称选择定则禁止，而 $2\orbitalP \to 1\orbitalS$ 可以通过电偶极跃迁快速衰变。微波一旦将原子从 $2\orbitalS$ 转入 $2\orbitalP$，它就会在到达探测器前衰变，存活的亚稳原子减少（图\ref{fig:lamb-readout}）。

他们固定微波频率，改变磁场强度，比较微波接通和关断时的探测电流，共振时电流减小得最多。在几个微波频率下测出共振磁场，再与磁子能级的理论曲线比较，可推得零场下 $2\orbitalS_{1/2}$ 高于 $2\orbitalP_{1/2}$，能量差对应的频率约为 $\qty{1}{GHz}$\cite{lamb}，即图\ref{fig:hydrogen-fine-lamb}最右列的间隔。
\begin{figure}[H]
    \centering\tikzfigure{lamb-readout}
    \caption{沿水平直箭头从制备端读到探测器，$N_{2\orbitalS}$ 为存活的亚稳原子数。两条波浪线分别标记入射微波 $\nu_{\mathrm{mw}}$ 和 $2\orbitalP \to 1\orbitalS$ 衰变发出的紫外光子 $\upgamma_{\mathrm{UV}}$。}
    \label{fig:lamb-readout}
\end{figure}

量子电动力学(quantum electrodynamics, QED)通过场量子化与重整化(renormalization)计算电子自能和真空极化对这一间隔的贡献。与精密谱线比较时，还须计入有限核大小和反冲。

\sectionstar{引力场中的量子干涉与量子引力}
\label{sec:qm-gravity}
中子干涉仪的两条臂若处在不同高度，地球引力就会移动干涉条纹。把近地面引力近似为均匀经典场，取竖直向上为 $z$ 轴，中子的势能为 $mgz$，其中 $g$ 是当地重力加速度。

在弱场近似下，沿未扰动的两条路径计算引力势能的一阶相位贡献。记两路在时刻 $t$ 的高度为 $z_a(t),z_b(t)$，从分束到合束有
\begin{equation}\label{eq:qm-gravity-phase}
    \Delta\varphi_g = -\frac{mg}{\hbar}\int_{t_i}^{t_f}
        [z_b(t) - z_a(t)]\,\d t .
\end{equation}
若中子在两条水平臂上各飞行 $T$，高度差恒为 $\Delta z$，且连接段的贡献相消，相位差便为 $-mg\Delta z\,T/\hbar$。
1975 年，\phantomsection\label{exp:qm-cow-gravity}Colella、Overhauser 和 Werner 用晶体将中子束分成两路，再使两路重新合束。他们倾斜整个干涉仪，改变两路的高度差，发现同一个出口的中子计数随倾角交替增减。相邻两个计数峰之间，两路的相对相位变化了 $2\pi$\cite{cow-neutrons}。
已知干涉仪的几何和中子的飞行时间，把当地测得的 $g$ 代入式\eqref{eq:qm-gravity-phase}，就能预言引力引起的条纹位移。

在这个中子实验的计算中，两路高度都是同一个实验室时间 $t$ 的函数。在给定时空背景上的量子场论中，度规同样预先给定，决定钟测得的固有时以及哪些事件能够互相发送信号。
广义相对论则把度规本身作为动力学变量，它会响应物质的能量和动量，还能以引力波传播。如果不同的量子分量对应不同的几何，就要指定物理钟和参照物，说明用哪些事件比较相位或先后次序。
这给量子引力带来了时间与可观测量的定义问题\cite{isham-time}。

在低能、弱曲率条件下，可在光滑背景上量子化小的引力涨落。把 Einstein 引力作为有效场论，按能量或曲率逐阶展开，每一阶只需有限个参数，就能算出长距离的量子修正\cite{donoghue-gravity}。
推向高能时，更高阶的曲率项会引入新的未知系数。能描述几何本身的量子态、恢复上述低能结果并在高能给出可检验预言的理论，仍有待实验判别。

\reading{精细结构的推导见 MIT 8.06 讲义 2\cite{mitfine}，微波测量见原论文\cite{lamb}。
    \hyperref[app:checks]{附录 A}联系谱线宽度与寿命，列出常用近似的适用尺度，并标明正文采用的定理与已知解。}
